% Generated by scripts/gen_blueprint.py from blueprint/nodes/12-not-found-in-the-sources.toml. Do not edit.

\chapter{Results not found in the sources read}
\label{chap:12-not-found-in-the-sources}

This chapter lists the library's results that were not found in the literature checked, sorted by how much they add. ``Not found'' means not found in about 110 sources (the papers and books of Diaz, Roy, Waldschmidt and Brownawell, and those they build on) nor in the works citing the key papers, as of 5 October 2026. It never means ``new'': each claim can fall to a source not yet read. ``Not routine'' means that the proof is not a substitution into a printed theorem. The source line of each statement says what is printed nearby. In the dependency graph, the results of the first list have a double border.

\textbf{Not found, and not routine.}
\begin{itemize}
\item Theorem~\ref{thm:circle_point_extension_two_by_three_configuration_iff}: the four-dimensional extensions $\Qbar+\Qbar u+\Qbar\bar u+\Qbar z$ of a point of a circle that carry a $2\times3$ configuration are exactly those containing $u^{2}$, $\bar u^{2}$ or some $1/(u-a)$. Its key step, Lemma~\ref{lem:two_by_three_configuration_forces_progression}, is the converse of a printed observation.
\item Theorem~\ref{thm:generic_circle_point_no_two_by_three_configuration}: no $2\times3$ configuration on generic data, whatever numbers are added. Its case $m=0$ is Roy's.
\item Theorem~\ref{thm:power_hull_strong_six_exp_configuration_iff}: in the power hulls $\Qbar u^{-k}+\Qbar u^{-1}+\Qbar+\Qbar u+\Qbar u^{k}$ a configuration exists only for $k=2,3$, so the strong six exponentials route at a candidate stops at $u^{3}$.
\item Theorems~\ref{thm:generic_quadratic_relation_is_norm} and~\ref{thm:generic_period_never_enters}: on generic data every quadratic relation is a multiple of the norm form, and the period never enters.
\item Theorems~\ref{thm:no_vanishing_coeff} and~\ref{thm:det_pencil_eq_conic}: a counterexample has no vanishing coefficient, and its only singular pencil is the norm form.
\end{itemize}

\textbf{Not found, but short: one substitution into a printed theorem, or a few lines.}
\begin{itemize}
\item Theorem~\ref{thm:circle_point_conjugate_pair_configuration_invisible_to_four_dimensional_extensions}: five-dimensional spaces $\Qbar+\Qbar u+\Qbar\bar u+\Qbar z+\Qbar\bar z$, with $z=u/(u^{2}-a)$ or $u/(u^{2}-a)^{2}$, carry a configuration that no four-dimensional extension of $\Qbar+\Qbar u+\Qbar\bar u$ inside them carries; and Proposition~\ref{prop:conj_stable_circle_point_extension_no_two_by_three_configuration}, the conjugation-stable case, which carries none. Their consequences at a candidate are printed (Corollary~\ref{cor:candidate_div_sq_sub_not_mem_logAlgTilde}).
\item Corollary~\ref{cor:dilog_half_irrational_or_exp_i_div_pi_transcendental}: $\mathrm{Li}_2(1/2)$ is irrational, or $e^{i\gamma/\pi}$ is transcendental for every rational $\gamma\neq0$; and its general form, Proposition~\ref{prop:recip_pi_log_of_rational_quadratic_relation}, which is Brownawell's Corollary~5 \cite{Bro74} with rational scaling.
\item Theorem~\ref{thm:diaz_2007_qr2_of_trdeg_one}: Diaz's (Qr2) in transcendence degree one; and Theorem~\ref{thm:log_pair_rigid_of_trdeg_one}, pair rigidity in transcendence degree one.
\item Corollaries~\ref{cor:torsion_rational_modulus_unique} and~\ref{cor:log_two_log_three_not_both_rational}: $(\log2)^{2}+\pi^{2}$ and $(\log3)^{2}+\pi^{2}$ are not both rational; and Corollary~\ref{cor:recip_pi_log_of_torsion_rational_modulus}: the two open boundary statements cannot both fail at rational data.
\item Theorem~\ref{thm:diaz_of_strong_five_exponentials}: the strong five exponentials conjecture implies Diaz's.
\end{itemize}

The companion note in the repository states and attributes all of these, with the passages read.
